% ===========================================================================
% MỞ ĐẦU
% ===========================================================================
\chapter*{MỞ ĐẦU}
\addcontentsline{toc}{chapter}{MỞ ĐẦU}
\markboth{MỞ ĐẦU}{MỞ ĐẦU}

\section*{1. Lý do chọn đề tài}
\addcontentsline{toc}{section}{1. Lý do chọn đề tài}

Trong những năm gần đây, chuyển đổi số trong lĩnh vực ngân hàng đã thúc đẩy mạnh
mẽ quá trình số hóa các sản phẩm và dịch vụ tín dụng. Cùng với sự phát triển của
ngân hàng số, nhu cầu cấp tín dụng nhanh chóng, thuận tiện và có khả năng xử lý
số lượng lớn hồ sơ ngày càng gia tăng. Đối với các tổ chức tín dụng, kiểm duyệt
tín dụng là một trong những nghiệp vụ quan trọng, có ảnh hưởng trực tiếp đến chất
lượng tài sản, khả năng kiểm soát rủi ro và hiệu quả kinh doanh. Do đó, việc ứng
dụng các mô hình học máy hỗ trợ cán bộ tín dụng trong quá trình đánh giá và ra
quyết định đang trở thành một yêu cầu thực tiễn.

Trong phương thức kiểm duyệt tín dụng truyền thống, thông tin khách hàng thường
được tổng hợp từ nhiều nguồn khác nhau như hồ sơ đăng ký vay, thông tin khách
hàng, lịch sử giao dịch, dữ liệu tín dụng và các thông tin liên quan khác. Quá
trình tổng hợp, kiểm tra và đánh giá dữ liệu có thể mất nhiều thời gian, đặc biệt
khi khối lượng hồ sơ lớn. Bên cạnh đó, việc đánh giá dựa nhiều vào các bộ quy tắc
nghiệp vụ cố định có thể hạn chế khả năng phát hiện các mối quan hệ phức tạp giữa
các đặc trưng của khách hàng và khả năng phát sinh rủi ro tín dụng.

Hạn chế nêu trên không nằm ở năng lực chuyên môn của cán bộ tín dụng mà nằm ở
\textit{bản chất của phương pháp đánh giá}. Một bộ quy tắc, dù được xây dựng công
phu đến đâu, cũng chỉ biểu diễn được quan hệ giữa từng chỉ tiêu riêng lẻ với
ngưỡng cho phép. Trong khi đó, rủi ro tín dụng thường phát sinh từ \textit{tổ
hợp} nhiều yếu tố: một khách hàng có thu nhập cao nhưng dòng tiền không ổn định,
thâm niên công việc ngắn và đang có quan hệ tín dụng với nhiều tổ chức cùng lúc
có thể rủi ro hơn một khách hàng thu nhập thấp hơn nhưng ổn định về mọi mặt. Bộ
quy tắc xét từng điều kiện độc lập khó nắm bắt được kiểu tương tác này.

Sự phát triển của các phương pháp học máy mở ra khả năng khai thác dữ liệu lịch
sử để xây dựng các mô hình dự đoán và hỗ trợ đánh giá rủi ro. Thay vì chỉ dựa
trên các điều kiện kiểm tra cố định, mô hình học máy có thể học các mối quan hệ
giữa đặc điểm khách hàng, hành vi tín dụng và kết quả tín dụng trong quá khứ để
đưa ra dự đoán đối với hồ sơ mới. Khi được tích hợp vào quy trình nghiệp vụ, kết
quả dự đoán có thể được sử dụng như một nguồn thông tin hỗ trợ cán bộ trong quá
trình kiểm duyệt tín dụng.

Tuy nhiên, việc ứng dụng học máy vào kiểm duyệt tín dụng không chỉ là bài toán
xây dựng mô hình dự đoán mà còn liên quan đến quá trình thu thập, chuẩn hóa, biến
đổi dữ liệu, xây dựng đặc trưng, tích hợp mô hình và cung cấp kết quả dự đoán một
cách ổn định. Thực tế cho thấy nhiều mô hình đạt chỉ tiêu tốt trong phòng thí
nghiệm nhưng không áp dụng được, vì thiếu các yếu tố mà môi trường vận hành đòi
hỏi: xác suất dự báo phải được hiệu chuẩn đúng để dùng làm căn cứ định giá, quyết
định phải giải trình được trước cán bộ tín dụng và cơ quan quản lý, và mô hình
phải xử lý được cả những hồ sơ thiếu dữ liệu. Vì vậy, cần có một giải pháp mang
tính hệ thống, trong đó mô hình học máy được kết hợp với quy trình xử lý dữ liệu
và các luật nghiệp vụ nhằm hỗ trợ quá trình kiểm duyệt tín dụng.

Xuất phát từ những yêu cầu trên, đề án lựa chọn đề tài
\textit{``\thesisTitle''} nhằm nghiên cứu, thiết kế và xây dựng thử nghiệm một mô
hình có khả năng tiếp nhận thông tin hồ sơ tín dụng, xử lý dữ liệu, tạo các đặc
trưng đầu vào, sử dụng mô hình học máy để đánh giá khả năng rủi ro và cung cấp
kết quả hỗ trợ kiểm duyệt. Đề tài có ý nghĩa cả về mặt khoa học và thực tiễn,
đồng thời có tính khả thi nhờ sự phát triển của các công nghệ xử lý dữ liệu, học
máy và các nền tảng triển khai hiện nay.

\section*{2. Tổng quan về vấn đề nghiên cứu}
\addcontentsline{toc}{section}{2. Tổng quan về vấn đề nghiên cứu}

Kiểm duyệt tín dụng là một bài toán đã được nghiên cứu rộng rãi trong lĩnh vực
tài chính -- ngân hàng. Các phương pháp truyền thống thường sử dụng chấm điểm tín
dụng và các bộ luật nghiệp vụ dựa trên những tiêu chí như thu nhập, lịch
sử tín dụng, dư nợ, khả năng trả nợ và đặc điểm của khoản vay. Những phương pháp
này có ưu điểm là dễ giải thích và phù hợp với các quy định nghiệp vụ, tuy nhiên
khả năng biểu diễn các mối quan hệ phức tạp trong dữ liệu còn hạn chế.

Cùng với sự gia tăng về quy mô và đa dạng của dữ liệu, nhiều nghiên cứu đã tập
trung áp dụng các thuật toán học máy cho bài toán đánh giá tín dụng. Một số
phương pháp phổ biến gồm hồi quy Logistic, Decision Tree, Random Forest, Support
Vector Machine, Gradient Boosting và các mô hình học sâu. Các nghiên cứu cho thấy
mô hình học máy có khả năng khai thác nhiều đặc trưng đồng thời và cải thiện khả
năng phân loại khách hàng có nguy cơ tín dụng cao hoặc thấp so với một số phương
pháp truyền thống.

Công trình đối sánh quy mô lớn của Lessmann và cộng sự~\cite{lessmann2015} trên
nhiều bộ dữ liệu tín dụng là một mốc tham chiếu quan trọng: các phương pháp học
tập hợp cho kết quả tốt hơn mô hình đơn lẻ, song mức chênh lệch so với hồi quy
logistic là \textit{không lớn} -- một kết luận thường bị bỏ qua khi lựa chọn mô
hình trong thực tiễn. Các tổng quan hệ thống gần đây~\cite{shi2022} tiếp tục
khẳng định xu hướng này, đồng thời chỉ ra rằng yếu tố quyết định chất lượng mô
hình thường nằm ở \textit{chất lượng và độ phong phú của đặc trưng} hơn là ở việc
lựa chọn thuật toán.

Bên cạnh vấn đề lựa chọn thuật toán, các nghiên cứu cũng tập trung vào các bước
xử lý dữ liệu như làm sạch dữ liệu, xử lý giá trị thiếu, lựa chọn đặc trưng, cân
bằng dữ liệu và đánh giá mô hình. Các chỉ số như Accuracy, Precision, Recall,
F1-score, ROC-AUC và Gini thường được sử dụng để đánh giá hiệu quả của mô hình
phân loại rủi ro tín dụng.

Tuy nhiên, một hạn chế thường gặp là nhiều nghiên cứu tập trung chủ yếu vào hiệu
quả dự đoán của mô hình, trong khi chưa giải quyết đầy đủ bài toán tích hợp mô
hình vào một quy trình hỗ trợ kiểm duyệt tín dụng hoàn chỉnh. Trong môi trường
thực tế, dữ liệu tín dụng có thể đến từ nhiều nguồn, có cấu trúc và chất lượng
khác nhau, đồng thời yêu cầu xử lý dữ liệu, kiểm tra điều kiện nghiệp vụ và cung
cấp kết quả trong thời gian phù hợp với quy trình kinh doanh.

Do đó, đề án tập trung nghiên cứu theo hướng kết hợp giữa xử lý dữ liệu, xây dựng
đặc trưng, mô hình học máy và luật nghiệp vụ để hình thành một mô hình hỗ trợ
kiểm duyệt tín dụng. Đề án không hướng tới việc thay thế hoàn toàn quyết định của
cán bộ hoặc tổ chức tín dụng mà tập trung xây dựng một công cụ hỗ trợ, cung cấp
kết quả đánh giá rủi ro và thông tin tham khảo nhằm nâng cao hiệu quả và tính
nhất quán của quá trình kiểm duyệt.

\section*{3. Mục tiêu nghiên cứu}
\addcontentsline{toc}{section}{3. Mục tiêu nghiên cứu}

\subsection*{3.1. Mục tiêu chung}

Đề án nhằm nghiên cứu, thiết kế và xây dựng thử nghiệm mô hình học máy hỗ trợ
kiểm duyệt tín dụng, trong đó khai thác dữ liệu khách hàng và thông tin tín dụng
để xây dựng các đặc trưng đầu vào, huấn luyện mô hình học máy và đưa ra kết quả
dự đoán mức độ rủi ro. Mô hình được xây dựng theo định hướng hỗ trợ quá trình
kiểm duyệt tín dụng, giúp cung cấp thêm cơ sở định lượng cho việc đánh giá hồ sơ
thay vì chỉ dựa trên các tiêu chí và quy tắc nghiệp vụ cố định.

\subsection*{3.2. Mục tiêu về mặt lý luận}

Về mặt lý luận, đề án tập trung hệ thống hóa và làm rõ cơ sở lý thuyết về kiểm
duyệt tín dụng, đánh giá rủi ro tín dụng và khả năng ứng dụng các phương pháp học
máy trong bài toán hỗ trợ ra quyết định tín dụng. Đề án nghiên cứu quy trình từ
thu thập và tiền xử lý dữ liệu, xây dựng đặc trưng, lựa chọn mô hình, huấn luyện,
đánh giá đến tích hợp mô hình. Trên cơ sở đó, đề án phân tích và
đánh giá mức độ phù hợp của các mô hình học máy đối với bài toán phân loại rủi ro
tín dụng thông qua các tiêu chí như Accuracy, Precision, Recall, F1-score và
ROC-AUC. Kết quả nghiên cứu góp phần làm rõ mối quan hệ giữa chất lượng dữ liệu,
đặc trưng đầu vào, lựa chọn mô hình và hiệu quả dự đoán trong bài toán tín dụng,
đồng thời bổ sung cơ sở tham khảo cho việc nghiên cứu ứng dụng học máy vào việc
hỗ trợ ra quyết định trong lĩnh vực tài chính -- ngân hàng.

\subsection*{3.3. Mục tiêu về mặt thực tiễn}

Về mặt thực tiễn, đề án hướng tới xây dựng một mô hình học máy thử nghiệm có khả
năng tiếp nhận thông tin hồ sơ tín dụng, thực hiện xử lý dữ liệu, tạo đặc trưng
để đánh giá rủi ro và trả về kết quả hỗ trợ kiểm duyệt. Việc ứng dụng mô hình học
máy giúp khai thác hiệu quả hơn dữ liệu lịch sử, hỗ trợ nhận diện các hồ sơ có
mức độ rủi ro khác nhau và cung cấp kết quả đánh giá một cách nhất quán. Qua đó,
có thể góp phần rút ngắn thời gian đánh giá hồ sơ, giảm khối lượng công việc thủ
công, nâng cao khả năng phát hiện rủi ro và hỗ trợ cán bộ tín dụng trong quá
trình ra quyết định. Kết quả của đề án cũng tạo cơ sở thực nghiệm để đánh giá khả
năng áp dụng mô hình học máy vào quy trình kiểm duyệt tín dụng và làm tiền đề cho
việc phát triển có quy mô lớn hơn trong môi trường thực tế.

\subsection*{3.4. Các mục tiêu cụ thể}

Từ hai định hướng trên, đề án đặt ra các mục tiêu cụ thể sau:

\begin{enumerate}[label=(\arabic*)]
  \item Hệ thống hóa cơ sở lý luận về tín dụng, kiểm duyệt tín dụng, rủi ro tín
  dụng và các phương pháp học máy áp dụng cho bài toán đánh giá tín dụng.

  \item Phân tích bài toán kiểm duyệt tín dụng: xác định các mức độ rủi ro cần dự
  báo, đặc trưng đầu vào, kết quả đầu ra và các yêu cầu kỹ thuật, pháp lý đối với
  mô hình.

  \item Xây dựng quy trình thu thập, tổng hợp, kiểm tra chất lượng và tiền xử lý
  dữ liệu hồ sơ tín dụng.

  \item Trích xuất và lựa chọn tập đặc trưng cho mô hình từ bốn nhóm nguồn dữ
  liệu: thông tin cá nhân và tài chính, tham số khoản vay, lịch sử tín dụng và
  quan hệ CIC, hành vi giao dịch tài khoản.

  \item Xây dựng nhãn rủi ro theo tiêu chuẩn nợ xấu và cửa sổ quan sát phù hợp;
  xác định điểm cắt phê duyệt theo các tiêu chí nghiệp vụ khác nhau.

  \item Thiết kế, huấn luyện, hiệu chuẩn và đánh giá so sánh các mô hình học máy
  theo bộ chỉ tiêu đầy đủ, bao gồm cả chỉ tiêu đặc thù của ngành tín dụng (KS,
  Gini, IV, PSI) và thời gian chấm điểm.

  \item Thiết kế và hiện thực bộ máy hỗ trợ ra quyết định kết hợp luật nghiệp vụ
  với mô hình học máy, có cơ chế giải thích kết quả cho từng hồ sơ.

  \item Mô phỏng luồng xử lý hồ sơ, thử nghiệm tích hợp qua dịch vụ REST và đối
  chiếu kết quả với quy trình thẩm định truyền thống chỉ dựa trên luật nghiệp vụ.
\end{enumerate}

\section*{4. Đối tượng và phạm vi nghiên cứu}
\addcontentsline{toc}{section}{4. Đối tượng và phạm vi nghiên cứu}

\subsection*{4.1. Đối tượng nghiên cứu}

Đối tượng nghiên cứu của đề án là mô hình học máy hỗ trợ kiểm duyệt tín dụng, tập
trung vào dữ liệu và quy trình đánh giá hồ sơ tín dụng. Cụ thể,
đề án nghiên cứu các dữ liệu liên quan đến khách hàng và khoản vay như thông tin
cá nhân, thông tin khoản vay, thu nhập, lịch sử tín dụng, tình trạng dư nợ và các
đặc trưng có liên quan đến khả năng phát sinh rủi ro tín dụng. Trên cơ sở các dữ
liệu này, đề án nghiên cứu phương pháp xử lý dữ liệu, xây dựng đặc trưng, huấn
luyện mô hình học máy và sử dụng kết quả dự đoán để hỗ trợ quá trình kiểm duyệt
tín dụng. Đối tượng nghiên cứu bao gồm mô hình học máy và kết quả đánh giá rủi ro
thành một quy trình hỗ trợ kiểm duyệt thống nhất.

\subsection*{4.2. Phạm vi nghiên cứu}

\textbf{Về nội dung:} Đề án giới hạn ở việc xây dựng và thử nghiệm mô hình hỗ trợ
đánh giá rủi ro tín dụng, bao gồm các bước chính: tiếp nhận dữ liệu hồ sơ,
tiền xử lý dữ liệu, xây dựng đặc trưng, áp dụng mô hình học máy, đánh giá kết quả
dự đoán và cung cấp kết quả hỗ trợ kiểm duyệt. Đề án không đi sâu nghiên cứu toàn
bộ quy trình nghiệp vụ tín dụng của ngân hàng, không xây dựng hệ thống quản lý
tín dụng hoàn chỉnh và không thay thế các quy định, chính sách hoặc quyết định
phê duyệt chính thức của tổ chức tín dụng. Các nội dung sau \textit{không} thuộc
phạm vi: định giá và thẩm định tài sản bảo đảm, quy trình giải ngân và thu hồi
nợ, mô hình ước lượng tổn thất khi vỡ nợ (LGD) và dư nợ tại thời điểm vỡ nợ
(EAD).

\textbf{Về không gian:} Nghiên cứu được thực hiện trong phạm vi bài toán kiểm
duyệt tín dụng tại tổ chức tín dụng/ngân hàng, tập trung vào việc xây dựng mô
hình và thử nghiệm trong môi trường nghiên cứu. Đề án không giới hạn vào
một địa bàn cụ thể và không thực hiện triển khai chính thức trên toàn bộ hệ thống
nghiệp vụ của một ngân hàng.

\textbf{Về dữ liệu:} Do đặc thù bảo mật của dữ liệu tín dụng cá nhân, đề án sử
dụng \textbf{dữ liệu mô phỏng} được sinh bằng chương trình tự xây dựng theo một
mô hình xác suất có cấu trúc. Thiết kế, cơ sở hiệu chỉnh và giới hạn của cách
tiếp cận này được trình bày minh bạch tại mục~\ref{sec:data-design}. Trong trường
hợp áp dụng trên dữ liệu thực tế, dữ liệu cần được ẩn danh và giới hạn quyền truy
cập nhằm bảo đảm yêu cầu về bảo mật thông tin.

\textbf{Về thời gian:} Đề án được thực hiện trong thời gian triển khai đề án tốt
nghiệp. Bộ dữ liệu sử dụng để huấn luyện và đánh giá mô hình được giới hạn trong
khoảng thời gian phù hợp với dữ liệu nghiên cứu và được phân chia thành các giai
đoạn phục vụ huấn luyện, kiểm thử và đánh giá mô hình. Các kết quả nghiên cứu
phản ánh phạm vi dữ liệu và điều kiện thực nghiệm trong thời gian thực hiện đề
án, chưa nhằm đánh giá khả năng vận hành của mô hình trong thời gian dài hoặc
trên quy mô dữ liệu sản xuất thực tế.

\textbf{Về kết quả:} Đề án tập trung đánh giá hai khía cạnh chính là hiệu quả của
mô hình học máy trong dự đoán rủi ro tín dụng và khả năng hỗ trợ quá trình kiểm
duyệt. Kết quả đầu ra của mô hình là thông tin đánh giá hoặc phân loại mức độ rủi
ro nhằm hỗ trợ cán bộ tín dụng và quy trình nghiệp vụ; quyết định
tín dụng cuối cùng vẫn thuộc về người hoặc bộ phận có thẩm quyền theo quy định
của tổ chức tín dụng.

\section*{5. Phương pháp nghiên cứu}
\addcontentsline{toc}{section}{5. Phương pháp nghiên cứu}

Đề án sử dụng kết hợp phương pháp nghiên cứu lý thuyết và phương pháp nghiên cứu
thực nghiệm nhằm nghiên cứu, xây dựng và đánh giá mô hình học máy trong hỗ trợ
kiểm duyệt tín dụng.

\textbf{Phương pháp nghiên cứu lý thuyết.} Đề án tiến hành thu thập, tổng hợp và
phân tích các tài liệu, giáo trình, công trình nghiên cứu và các tài liệu chuyên
ngành liên quan đến tín dụng, kiểm duyệt tín dụng, đánh giá rủi ro tín dụng, chấm
điểm tín dụng và ứng dụng học máy trong lĩnh vực tài chính -- ngân hàng. Trên cơ
sở nghiên cứu lý thuyết, đề án làm rõ các khái niệm, đặc điểm của bài toán đánh
giá tín dụng, các yếu tố ảnh hưởng đến rủi ro tín dụng và nguyên lý hoạt động của
các mô hình học máy phù hợp. Đồng thời, nghiên cứu các phương pháp tiền xử lý dữ
liệu, lựa chọn và xây dựng đặc trưng, huấn luyện và đánh giá mô hình. Kết quả
nghiên cứu lý thuyết được sử dụng làm cơ sở để lựa chọn mô hình học máy, thiết kế
quy trình xử lý dữ liệu và xây dựng mô hình học máy hỗ trợ kiểm duyệt tín dụng.

\textbf{Phương pháp nghiên cứu thực nghiệm.} Đề án tiến hành xây dựng môi trường
thực nghiệm để kiểm chứng các nội dung nghiên cứu lý thuyết trên bộ dữ liệu tín
dụng phục vụ nghiên cứu. Dữ liệu được phân tích và xử lý thông qua các bước làm
sạch, xử lý dữ liệu thiếu, mã hóa dữ liệu, xây dựng và lựa chọn đặc trưng trước
khi đưa vào mô hình. Một số mô hình học máy phù hợp với bài toán phân loại rủi ro
tín dụng được xây dựng, huấn luyện và kiểm thử trên dữ liệu thực nghiệm. Cụ thể,
phương pháp thực nghiệm được triển khai qua bốn nhóm công việc:

\begin{itemize}
  \item \textbf{Phân tích dữ liệu:} thống kê mô tả tập dữ liệu, kiểm tra tính hợp
  lệ theo ràng buộc nghiệp vụ, phân tích giá trị khuyết thiếu và giá trị ngoại
  lệ; đánh giá năng lực dự báo của từng đặc trưng thông qua chỉ số Information
  Value.

  \item \textbf{Thực nghiệm học máy:} huấn luyện, hiệu chuẩn và so sánh các mô
  hình theo phương pháp kiểm định ngoài thời gian (\eng{out-of-time validation});
  tinh chỉnh siêu tham số bằng kiểm định chéo theo chuỗi thời gian; giải thích mô
  hình bằng phương pháp SHAP và phép hoán vị.

  \item \textbf{Mô phỏng quy trình kiểm duyệt:} xây dựng bộ máy hỗ trợ ra quyết
  định, mô phỏng luồng xử lý hồ sơ, xác định ngưỡng phê duyệt theo ba tiêu chí
  nghiệp vụ và thử nghiệm tích hợp qua dịch vụ REST.

  \item \textbf{Đánh giá đối chiếu:} so sánh kết quả của mô hình đề xuất với quy
  trình thẩm định truyền thống chỉ dựa trên luật nghiệp vụ, và đối chiếu với các
  công trình nghiên cứu liên quan.
\end{itemize}

\section*{6. Kết cấu của đề án}
\addcontentsline{toc}{section}{6. Kết cấu của đề án}

Ngoài phần Mở đầu, Kết luận, Danh mục tài liệu tham khảo và Phụ lục, nội dung đề
án được trình bày trong ba chương:

\begin{itemize}
  \item \textbf{Chương 1 -- Cơ sở lý luận:} trình bày tổng quan về tín dụng và
  kiểm duyệt tín dụng, tổng quan về học máy trong đánh giá tín dụng, dữ liệu và
  đặc trưng trong bài toán tín dụng, các mô hình học máy áp dụng trong đánh giá
  tín dụng cùng cơ sở so sánh và lựa chọn mô hình phù hợp.

  \item \textbf{Chương 2 -- Xây dựng mô hình học máy trong hỗ trợ kiểm duyệt tín
  dụng:} phân tích bài toán kiểm duyệt tín dụng; trình bày quy trình thu thập
  và tổng hợp dữ liệu, tiền xử lý và làm sạch dữ liệu, trích xuất và lựa chọn đặc trưng, xây
  dựng nhãn và xác định quy tắc phê duyệt; thiết kế mô hình học máy và quy trình
  huấn luyện, tối ưu mô hình.

  \item \textbf{Chương 3 -- Triển khai thực nghiệm và đánh giá kết quả:} trình
  bày kịch bản và môi trường thực nghiệm, kết quả huấn luyện các mô hình, đánh
  giá hiệu năng và so sánh giữa các thuật toán, mô phỏng và giải thích mô hình;
  cuối cùng là phần đánh giá chung, đối chiếu với quy trình thẩm định
  truyền thống và với các công trình nghiên cứu liên quan.
\end{itemize}
